\subsection{阴毛修剪、脱毛与私处外观护理（安全与误区）}

\begin{itemize}
	\item \textbf{阴毛的生理功能}：阴毛可减少皮肤摩擦、缓冲机械刺激、帮助散发外阴的汗液与气味，是身体正常的生理结构。是否保留、如何修整，属于个人审美与卫生习惯的选择，\textbf{没有"必须剃掉"的医学理由}。
	\item \textbf{修剪 vs.\ 完全剃除}：修剪（保留 0.5--1 cm）相对安全；完全剃除后新生毛发短硬、回刺入皮，容易引起\textbf{假性毛囊炎}（红肿、瘙痒、小脓疱）与毛发向内生长（内生毛），皮肤屏障受损还会增加感染与性传播感染的风险。
	\item \textbf{比基尼蜡与激光脱毛}：蜡脱有效但疼痛明显，可能引起毛囊炎、色素沉着；激光/强脉冲光脱毛效果较持久，应选择正规医疗机构操作，治疗前避免暴晒、期间避免使用刺激性护肤品，孕期与光敏感者慎用。
	\item \textbf{剃毛的安全要点}：如需剃除，用\textbf{专用剃刀}（勿与他人共用）、温水软化毛发、顺着毛发生长方向刮、使用剃须泡或润滑、事后涂抹温和保湿品，避免干刮与反复刮同一部位；不要用镊子拔毛（易感染与内生毛）。
	\item \textbf{私处"美白""漂白"产品的风险}：市售所谓"嫩红""私处美白"产品常含\textbf{氢醌、汞、糖皮质激素或强酸}成分，外阴皮肤薄而敏感，使用后可能引起\textbf{化学性灼伤、皮炎、色素紊乱甚至全身吸收中毒}，切勿轻信。外阴颜色深浅与肤色、激素水平、年龄、摩擦有关，\textbf{个体差异属于正常}。
	\item \textbf{私处香氛/香水}：外阴专用香氛、带香味的喷雾与湿巾，其香料与酒精成分\textbf{破坏阴道酸性环境、诱发接触性皮炎}，有异味时更应就医查因，而不是"用香味盖住"。
	\item \textbf{外观焦虑的正名}：小阴唇大小、形态、颜色以及是否对称，个体差异极大，绝大多数都属于正常变异。若因外观而长期焦虑、影响性生活或日常（如骑车、行走时疼痛），可到正规医院妇科或整形科咨询，\textbf{但应避免受网络营销影响而盲目手术}。
	\item \textbf{异常信号}：外阴持续性瘙痒、疼痛、溃疡、增厚变白、疣状赘生物、反复毛囊炎或出血，应及时就诊，不要自行涂抹激素或抗生素药膏。
\end{itemize}
